\problem{Newton interpolation and Runge's phenomenon}
    Consider the \textit{Runge} function
    \begin{equation}
        f_A(x) = \frac{1}{1+25x^2}
    \end{equation}
    as well as
    \begin{equation}
        f_B(x) = \sin(x)
    \end{equation}
    and
    \begin{itemize}
        \item equidistant evaluation points
            \begin{equation}
                x_i = \frac{2i}{m} - 1, \quad i = 0,...,m
            \end{equation}
        \item as well as \textit{Chebyshev-nodes}
            \begin{equation}
                \tilde{x}_i=\cos \left(\frac{(2 i+1) \pi}{2(m+1)}\right), \quad i=0, \ldots, m
            \end{equation}
    \end{itemize}
    In this task, we will perform \textbf{Newton interpolation}

    \begin{equation}
        p(x) = \sum_{j=0}^{m} d_j n_j(x), \quad n_j(x) = \prod_{i = 0}^{j - 1} (x-x_i)
    \end{equation}

    with the coefficients given by the finite differences

    \begin{equation}
        d_j = [x_0,...,x_j] f
    \end{equation}

    where for $j<i$ we recursively define

    \begin{equation}
        \begin{aligned}
            \left[x_j\right]f & = f_j, \\
            \left[x_j, \ldots, x_i\right] f &=\frac{\left[x_{j+1}, \ldots, x_i\right] f-\left[x_j, \ldots, x_{i-1}\right] f}{x_i-x_j}
        \end{aligned}
    \end{equation}

    \begin{enumerate}
        \item Implement Newton divided differences in Python.
        \item Implement the modified Horner scheme to evaluate $p(x)$ via
            \begin{equation}
                z_m = d_m, z_j = d_j + (x - x_j) z_{j+1} \quad \rightarrow \quad z_0 = p(x)
            \end{equation}
        \item For both $f_A$ and $f_B$ perform Newton interpolation for $m = 3, 9, 24, 48$ both
        with equidistant and Chebyshev nodes. Plot the results. For plotting $p(x)$ use
        $m = 200$ evaluation points. Interpret your results.
    \end{enumerate}

\begin{solutionblock}
\part
Todo
\end{solutionblock}

\pagebreak